% !Mode:: "TeX:UTF-8"

\def\usewhat{xelatex}                                % 使用现代 XeTeX 兼容编译方式
                                                     % 优先使用 macOS 中文字体，其他平台自动回退 Fandol。
\documentclass[12pt,openany,oneside,fontset=none]{ctexbook}
                                                     % 本科生毕业论文通常采用单页排版
\IfFontExistsTF{STSong}{
  \setCJKmainfont[BoldFont=STHeiti]{STSong}
  \setCJKsansfont[AutoFakeBold=2.5]{STHeiti}
  \setCJKmonofont{STFangsong}
  \setCJKfamilyfont{zhsong}[BoldFont=STHeiti]{STSong}
  \setCJKfamilyfont{zhhei}[AutoFakeBold=2.5]{STHeiti}
  \setCJKfamilyfont{zhkai}{STKaiti}
  \setCJKfamilyfont{zhfs}{STFangsong}
}{
  \setCJKmainfont[BoldFont=FandolSong-Bold]{FandolSong-Regular}
  \setCJKsansfont[BoldFont=FandolHei-Bold]{FandolHei-Regular}
  \setCJKmonofont{FandolFang-Regular}
  \setCJKfamilyfont{zhsong}[BoldFont=FandolSong-Bold]{FandolSong-Regular}
  \setCJKfamilyfont{zhhei}[BoldFont=FandolHei-Bold]{FandolHei-Regular}
  \setCJKfamilyfont{zhkai}{FandolKai-Regular}
  \setCJKfamilyfont{zhfs}{FandolFang-Regular}
}
\input{setup/package}                                % 定义本文所使用宏包
\graphicspath{{figures/}{figures/srs/}}               % 报告公共图片与 SRS 图片目录
\begin{document}                                     % 开始全文
	\input{setup/format}                             % 完成对论文各个部分格式的设置
	% !Mode:: "TeX:UTF-8"


%%%%%%%%%%%%%%%%%%%%%%%%%%%%%%%%%%%%%%%%%%%%%%%%%%%%%%%%%%%%%%%
%%  可通过对 setup/format.tex中                               %%
%%  第243行 \setlength{\@title@width}{5cm}中 5cm 这个参数来   %%
%%  控制封面中下划线的长度。                                   %%
%%%%%%%%%%%%%%%%%%%%%%%%%%%%%%%%%%%%%%%%%%%%%%%%%%%%%%%%%%%%%%

% 交叉验收时将下行改为 \srsanonymoustrue，即可生成匿名封面与 PDF 元数据。
\newif\ifsrsanonymous
\srsanonymousfalse

\cheading{软件工程综合实践~\textperiodcentered~需求分析附录~1}      % 正文页眉
\ccovertitle{《软件工程综合实践》结课报告}                       % 封面标题

%%%%%%%%%%%%%%%%%%%%%%%%%%%%%%%%%%%%%%%%%%%%%%%%%%%%%%%%%%%%%
%%%%%%%%%% 以下为论文的基本信息，需要由作者进行修改 %%%%%%%%%%%%
%%%%%%%%%%%%%%%%%%%%%%%%%%%%%%%%%%%%%%%%%%%%%%%%%%%%%%%%%%%%%
\ctitle{轻量级外卖服务平台的需求规格说明书}    % 封面用论文标题，自己可手动断行
\caffil{软件学院}       % 学院名称（已与页眉统一）
\csubject{软件工程}     % 专业名称
\cgrade{2024级}           % 年级
\ifsrsanonymous
  \cauthor{匿名}
  \cstuid{\xiaowu 匿名}
\else
  \cauthor{范文丽、李承文、赵紫怡、仲禹潼}          % 学生姓名（按组员登记顺序）
  \cstuid{\xiaowu 3024244438、3024232129、3024244137、3024244140}     % 学号，与姓名顺序一致
\fi

\ifsrsanonymous
  \hypersetup{pdftitle={轻量级外卖服务平台的需求规格说明书},pdfauthor={匿名},pdfsubject={软件工程综合实践},pdfkeywords={软件需求规格说明书, 外卖系统, 骑手配送}}
\else
  \hypersetup{pdftitle={轻量级外卖服务平台的需求规格说明书},pdfauthor={范文丽、李承文、赵紫怡、仲禹潼},pdfsubject={软件工程综合实践},pdfkeywords={软件需求规格说明书, 外卖系统, 骑手配送}}
\fi

\cdate{\the\year~年~\the\month~月~\the\day~日}  % 论文完成日期，不需要修改，自动生成

	\frontmatter                                     % 以下是论文导言部分，包括论文的封面，中英文摘要和中文目录
	\hypersetup{pageanchor=false}                    % 前置页不创建重复的数字页锚点
	\fancypagestyle{plain}{							 % 正文前均无页眉
		\fancyhf{}
		\renewcommand{\headrulewidth}{0 pt}
		\fancyfoot[C]{\song\xiaowu~\thepage~}
	}

	\makecover 			% 封面

	\input{setup/toc} % 目录

	\mainmatter\hypersetup{pageanchor=true}\defaultfont\sloppy\raggedbottom
	\makeatletter
	\fancypagestyle{plain}{                              % 设置正文眉页脚风格
		\fancyhf{}
		\fancyhead[C]{\song\wuhao \@cheading}            % 页眉格式
		\fancyfoot[C]{\song\xiaowu ~\thepage~}           % 页脚格式
		\renewcommand{\headrulewidth}{0.5pt}
		\renewcommand{\footrulewidth}{0pt}
	}
	\makeatother
	\setcounter{page}{1}                                 % 单独从 1 开始编页码
	\titleformat{\chapter}{\centering\xiaosan\hei}{\chaptername}{2em}{} % 恢复chapter标题格式要求

	%%%%%% 正文入口：章节内部再按主题拆分，便于并行协作 %%%%%%
	% SRS 组装入口。请在各子文件中协作，避免多人同时修改本文件。
\chapter{软件需求规格说明书}

\section{引言}

\subsection{文档目的}

本需求规格说明书定义“轻量级外卖服务平台”的业务边界、功能行为、业务规则、外部接口和质量要求，作为需求评审、架构设计、测试驱动开发、前后端联调和交叉验收的共同基线。文档面向课程教师、验收小组、项目经理、开发人员和测试人员。

\subsection{项目背景}

项目围绕外卖业务的四方协同场景展开。顾客需要完成浏览、选购、结算和收货；商家需要维护店铺、商品并履行订单；骑手需要承接配送任务；管理员需要完成准入与异常治理。系统在完成上述基础闭环的同时，引入 AI 客服与智能点餐、营销资产和数据看板，但不让任何外部智能服务成为基础下单和配送的单点依赖。

\subsection{产品愿景与目标}

\begin{enumerate}
  \item 建立顾客、商家、骑手和管理员职责清晰的外卖履约闭环。
  \item 以订单和配送任务的独立状态流转保证业务可理解、可追踪和可测试。
  \item 以权限受控的 AI 客服连接菜单搜索、订单查询和智能点餐，展示智能服务与传统交易系统的可控集成。
  \item 对库存、订单、营销资产和骑手接单提供后端业务校验、幂等和并发保护。
  \item 使需求、接口、数据对象、自动化测试和验收结果保持可追踪关系。
\end{enumerate}

\subsection{优先级定义}

\begin{longtable}{|L{0.16\linewidth}|L{0.30\linewidth}|L{0.44\linewidth}|}
\caption{需求优先级定义}\label{tab:priority}\\ \hline
\textbf{级别} & \textbf{含义} & \textbf{验收处理} \\ \hline
P0 & 核心基线 & 必须实现并通过正常、异常和边界验收。 \\ \hline
P1 & 计划完成的亮点 & 属于本版本基线，应实现并演示完整闭环。 \\ \hline
P2 & 可选增强 & 不作为基线验收失败条件；启用时仍须符合权限、安全和降级要求。 \\ \hline
\end{longtable}

\subsection{术语与缩略语}

\begin{longtable}{|L{0.22\linewidth}|L{0.68\linewidth}|}
\caption{术语与缩略语}\label{tab:terms}\\ \hline
\textbf{术语} & \textbf{定义} \\ \hline
有效订单 & 已支付且未取消的订单；营业额与销量默认仅统计已完成的有效订单。 \\ \hline
配送任务 & 订单进入配送阶段后生成的独立履约记录，记录骑手、任务状态和关键时间。 \\ \hline
活动任务 & 处于待接单、已接单、已取餐、配送中或异常待处理状态的配送任务。 \\ \hline
资产流水 & 钱包、积分和优惠券因业务动作产生的不可直接覆盖的变更记录。 \\ \hline
幂等 & 同一业务请求重复提交时，系统不产生重复扣款、扣库存、返还或状态跳转。 \\ \hline
AI 工具调用 & AI 客服在用户授权与后端权限校验后，对菜单搜索、订单查询等受控接口的调用。 \\ \hline
降级 & 外部 AI、地图或媒体服务不可用时，返回明确提示并切换到普通搜索、文本地址或手工操作。 \\ \hline
\end{longtable}

\subsection{参考资料}

\begin{longtable}{|L{0.34\linewidth}|L{0.18\linewidth}|L{0.38\linewidth}|}
\caption{参考资料}\label{tab:references}\\ \hline
\textbf{资料} & \textbf{提供方} & \textbf{用途} \\ \hline
《软件工程综合实践》项目要求 & 课程组 & 确定前后端分离、测试先行、RESTful 接口、需求变更和交叉验收等过程要求。 \\ \hline
轻量级外卖服务平台公开仓库 & 项目组 & 保存对外源代码、需求、接口、测试和验收材料。 \\ \hline
\end{longtable}

\section{项目总体说明}

\subsection{涉众分析}

\begin{longtable}{|L{0.14\linewidth}|L{0.30\linewidth}|L{0.46\linewidth}|}
\caption{涉众分析}\label{tab:stakeholders}\\ \hline
\textbf{涉众} & \textbf{主要目标} & \textbf{待解决问题} \\ \hline
顾客 & 高效、透明地完成点餐 & 快速查找合适菜品，理解优惠与应付金额，查看履约进度，处理取消、收货和评价。 \\ \hline
商家 & 维护店铺并稳定履约 & 控制营业状态、商品和库存，按状态处理订单，与骑手完成取餐交接。 \\ \hline
骑手 & 安全承接并完成配送 & 查看可接任务，避免重复抢单，按次序完成取餐、配送、送达和异常上报。 \\ \hline
管理员 & 保证平台准入和运行秩序 & 审核商家、店铺和骑手，处理账号、评价和配送异常，查看平台统计。 \\ \hline
验收人员 & 根据稳定基线复现结果 & 在统一环境和初始数据下，按需求编号验证正常、异常、权限和边界场景。 \\ \hline
\end{longtable}

\subsection{角色与权限边界}

\begin{longtable}{|L{0.14\linewidth}|L{0.28\linewidth}|L{0.28\linewidth}|L{0.20\linewidth}|}
\caption{角色权限边界}\label{tab:role-permission}\\ \hline
\textbf{角色} & \textbf{可读数据} & \textbf{可写数据} & \textbf{禁止行为} \\ \hline
顾客 & 公开店铺商品、本人订单与资产 & 本人资料、地址、购物车、订单和评价 & 访问他人订单、直接改价或跳过订单状态 \\ \hline
商家 & 本人店铺、商品、订单和统计 & 本人店铺资源、合法的商家履约状态 & 操作非本店数据、修改顾客原评价 \\ \hline
骑手 & 可接任务和本人活动/历史任务 & 本人任务的到店、取餐、送达与异常信息 & 查看非必要顾客资料、操作他人任务 \\ \hline
管理员 & 审核材料、平台治理数据和汇总统计 & 审核、账号状态、评价展示状态和异常处理结果 & 直接覆盖资产流水或删除审计记录 \\ \hline
\end{longtable}

\subsection{系统范围与边界}

\srsfigure{figures/srs/01-system-context.pdf}{系统上下文图}{fig:system-context}

图~\ref{fig:system-context} 应展示四类登录角色、轻量级外卖服务平台以及 AI、地图和对象存储等外部依赖。外部服务仅提供候选结果或导航能力，不直接写入订单、库存和用户资产。

\subsection{功能范围}

\begin{longtable}{|L{0.12\linewidth}|L{0.24\linewidth}|L{0.54\linewidth}|}
\caption{分级功能范围}\label{tab:scope}\\ \hline
\textbf{级别} & \textbf{模块} & \textbf{范围} \\ \hline
P0 & 核心交易 & 账号、地址、店铺商品、购物车、订单、模拟支付、取消、查询与确认收货。 \\ \hline
P0 & 履约和治理 & 商家接单/制作/出餐，骑手申请/上线/接单/取餐/配送/异常，管理员准入审核与账号治理。 \\ \hline
P1 & 智能体验 & 权限受控的 AI 客服、订单查询、菜单搜索和结构化智能点餐。 \\ \hline
P1 & 增强业务 & 模拟钱包、积分、单张优惠券、固定会员权益、评价闭环、三类数据看板、地图导航入口和实时状态通知。 \\ \hline
P2 & 可选增强 & 图片识菜、语音输入、每日幸运主题、AI 商家经营摘要和实时 GPS 轨迹。 \\ \hline
\end{longtable}

\subsection{范围外事项}

本期不接入真实微信/支付宝资金清算，不实现复杂自动派单、骑手薪资与劳动关系、人工客服坐席、跨平台物流、任意组合的促销叠加和持续实时 GPS 轨迹。地图仅提供取餐和送餐导航入口，文本地址始终可用。

\subsection{运行环境与假设}

系统采用 Vue 3 前端、Spring Boot 后端和 MySQL 数据库。基线验收使用项目提供的演示账号与初始化数据，只要浏览器能访问项目服务即可完成基础业务。AI、地图和对象存储通过适配接口调用，其凭证只保存在后端环境变量中。

\section{业务模型与规则}

\subsection{核心业务对象}

\begin{longtable}{|L{0.19\linewidth}|L{0.25\linewidth}|L{0.46\linewidth}|}
\caption{核心业务对象}\label{tab:domain-objects}\\ \hline
\textbf{对象} & \textbf{责任} & \textbf{关键约束} \\ \hline
User / Role & 保存账号、个人资料和角色 & 同一用户可具有多个角色；禁用账号不得登录。 \\ \hline
Store / Product & 表示店铺、分类、商品和库存 & 未通过审核或休息店铺不可接收新订单；下架或库存为零的商品不可购买。 \\ \hline
Cart / CartItem & 保存顾客当前待结算商品 & 一个购物车只包含同一店铺商品；金额仅供展示，提交订单时重新计算。 \\ \hline
Order / OrderItem & 保存交易状态、金额和商品/地址快照 & 金额由后端计算；历史订单不因商品或地址修改而改变。 \\ \hline
Payment\-Record & 记录模拟支付或钱包支付结果 & 支付状态与订单状态分离；同一订单仅有一次有效支付成功。 \\ \hline
DeliveryTask & 保存履约状态、骑手和关键时间 & 一个订单可有多条历史配送尝试，但同一时刻最多一条活动任务。 \\ \hline
Review & 保存评分、文字、图片和商家回复 & 仅已完成的本人订单可评价；每个订单最多一条有效顾客评价。 \\ \hline
\end{longtable}

\srsfigure{figures/srs/02-core-domain-class.pdf}{核心交易与履约领域类图}{fig:core-domain-class}

\subsection{营销资产与个性化对象}

钱包、积分和优惠券均通过账户与流水追踪变化；会员保存方案、生效时间和到期时间；用户偏好仅影响展示、搜索和推荐，不改变历史订单。

\srsfigure{figures/srs/03-assets-preference-class.pdf}{用户资产与个性化领域类图}{fig:assets-class}

\subsection{外卖履约主流程}

\srsfigure{figures/srs/04-fulfillment-activity.pdf}{外卖订单履约活动图}{fig:fulfillment-activity}

图~\ref{fig:fulfillment-activity} 使用顾客、订单服务、商家、配送服务、骑手和通知服务六个泳道，必须包含库存不足、支付失败、商家拒单、用户取消和配送异常分支。

\subsection{订单状态规则}

\srsfigure{figures/srs/05-order-state.pdf}{订单状态机图}{fig:order-state}

\begin{longtable}{|L{0.16\linewidth}|L{0.18\linewidth}|L{0.18\linewidth}|L{0.38\linewidth}|}
\caption{订单状态转换}\label{tab:order-state}\\ \hline
\textbf{当前状态} & \textbf{事件/执行者} & \textbf{目标状态} & \textbf{条件与副作用} \\ \hline
无 & 提交订单/顾客 & 待支付 & 校验店铺、商品、地址和优惠；生成快照并预占库存与优惠券。 \\ \hline
待支付 & 支付成功/顾客 & 待商家接单 & 确认支付流水和资产核销，向商家发送新订单通知。 \\ \hline
待支付 & 主动取消或 15 分钟超时/顾客或系统 & 已取消 & 释放预占库存和优惠券，不生成扣款。 \\ \hline
待商家接单 & 接单/商家 & 制作中 & 校验订单属于本店且未取消，记录接单时间。 \\ \hline
待商家接单 & 拒单或顾客取消 & 已取消 & 记录取消方和原因，恢复库存并原路返还资产。 \\ \hline
制作中 & 确认出餐/商家 & 待骑手接单 & 生成一条待接配送任务；顾客不再可自助取消。 \\ \hline
待骑手接单 & 接单/骑手 & 待取餐 & 通过带状态条件的原子更新绑定唯一骑手。 \\ \hline
待取餐 & 确认取餐/骑手 & 配送中 & 仅本任务骑手可操作，记录取餐时间。 \\ \hline
配送中 & 确认送达/骑手 & 已送达 & 记录送达时间，向顾客发送确认收货提示。 \\ \hline
已送达 & 确认收货或 24 小时超时/顾客或系统 & 已完成 & 纳入营业额和销量统计，发放积分并开放评价入口。 \\ \hline
\end{longtable}

\subsection{配送任务状态规则}

\srsfigure{figures/srs/06-delivery-state.pdf}{配送任务状态机图}{fig:delivery-state}

\begin{longtable}{|L{0.18\linewidth}|L{0.18\linewidth}|L{0.20\linewidth}|L{0.34\linewidth}|}
\caption{配送任务状态转换}\label{tab:delivery-state}\\ \hline
\textbf{当前状态} & \textbf{操作} & \textbf{目标状态} & \textbf{约束} \\ \hline
待接单 & 骑手接单 & 已接单 & 骑手已审核、在线且无超出上限的活动任务；并发时仅一人成功。 \\ \hline
已接单 & 确认到店 & 已接单 & 记录到店时间，不额外增加状态。 \\ \hline
已接单 & 确认取餐 & 已取餐 & 商家已出餐，操作人为当前绑定骑手。 \\ \hline
已取餐 & 开始配送 & 配送中 & 展示必要收货信息和地图导航入口。 \\ \hline
配送中 & 确认送达 & 已送达 & 记录送达时间，完成后限制骑手继续读取敏感地址。 \\ \hline
已接单/已取餐/配送中 & 上报异常 & 异常待处理 & 必须选择异常类型并填写说明，不允许骑手直接取消订单。 \\ \hline
异常待处理 & 恢复/改派/终止 & 原状态/已取消 & 由管理员记录处理结果；改派时终止原任务并建立新任务。 \\ \hline
\end{longtable}

\subsection{订单与配送联动}

\begin{longtable}{|L{0.20\linewidth}|L{0.22\linewidth}|L{0.20\linewidth}|L{0.28\linewidth}|}
\caption{订单与配送联动矩阵}\label{tab:order-delivery-map}\\ \hline
\textbf{订单状态} & \textbf{配送状态} & \textbf{可操作角色} & \textbf{主要操作} \\ \hline
待支付/待商家接单/制作中 & 无 & 顾客/商家 & 支付、取消、接单、拒单、出餐 \\ \hline
待骑手接单 & 待接单 & 骑手 & 查看任务、接单 \\ \hline
待取餐 & 已接单 & 骑手/商家 & 到店、核对、取餐 \\ \hline
配送中 & 已取餐/配送中 & 骑手 & 导航、送达、上报异常 \\ \hline
已送达/已完成 & 已送达 & 顾客 & 确认收货、评价 \\ \hline
已取消 & 无/已取消 & 系统/管理员 & 释放资源、保存取消和处理记录 \\ \hline
\end{longtable}

\subsection{资产结算与回退}

\srsfigure{figures/srs/07-settlement-sequence.pdf}{订单结算与资产处理顺序图}{fig:settlement-sequence}

结算顺序固定为“商品小计→会员折扣→优惠券→积分抵扣→钱包/模拟支付”。每单最多使用一张优惠券；优惠后应付金额不得小于零；订单取消时按原流水生成反向流水，不直接删除原记录。

\subsection{AI 客服与智能点餐流程}

\srsfigure{figures/srs/08-ai-assistant-sequence.pdf}{AI 客服与智能点餐顺序图}{fig:ai-assistant}

AI 客服可回答平台规则，也可在登录且通过后端资源归属校验后查询当前用户的订单状态。智能点餐根据预算、口味、忌口和当前可售菜单返回结构化候选商品；候选结果必须经商品服务二次校验，并由顾客明确确认后才能加入购物车。AI 不得直接支付、取消订单、修改库存或调整用户资产。

\section{功能性需求}

\subsection{功能模块总览}

\srsfigure{figures/srs/09-feature-map.pdf}{系统功能模块图}{fig:feature-map}

系统功能分为公共账号与通知、顾客端、商家端、骑手端、管理端、智能服务、营销资产与个性化七个模块。每项需求使用稳定编号，后续接口、数据库、测试和验收材料必须引用同一编号。

\subsection{系统总用例}

\srsfigure{figures/srs/10-system-use-case.pdf}{系统总用例图}{fig:system-use-case}

总用例图只展示四类业务角色、游客和外部服务与一级用例的关系，各角色内部用例见后续分图和用例表。

\subsection{用例编写规则}

每个用例包含执行者、触发条件、前置条件、基本流程、异常与边界、后置条件、业务规则和可直接执行的验收标准。需求中的“顺序”、“仅本人”、“不重复”等限制均应由后端再次校验，不得只依赖前端隐藏按钮。

\input{body/srs/functions/common}
\input{body/srs/functions/customer}
\input{body/srs/functions/merchant}
\input{body/srs/functions/rider}
\input{body/srs/functions/admin}
\input{body/srs/functions/intelligent_assets}
\section{外部接口与界面需求}

\subsection{四端页面结构}

\srsfigure{figures/srs/15-page-information-architecture.pdf}{四端页面信息架构图}{fig:page-ia}

顾客端以店铺浏览、购物车、订单和个人中心为主导航；商家端以店铺、商品、订单和看板为主；骑手端以可接任务、当前任务和历史为主；管理端以审核、治理和平台看板为主。角色切换后必须进入相应工作台，不保留上一角色的私有页面状态。

\subsection{关键页面原型}

\srsfigure{figures/srs/16-key-page-wireframes.pdf}{关键页面低保真原型组}{fig:key-wireframes}

原型组至少包含顾客 AI 点餐页、订单履约详情页、商家订单工作台、骑手当前任务页和管理员审核/看板页。原型用于约定信息层级和状态反馈，视觉色彩可在不改变需求的前提下调整。

\subsection{页面通用状态}

\begin{longtable}{|L{0.17\linewidth}|L{0.31\linewidth}|L{0.42\linewidth}|}
\caption{页面通用状态}\label{tab:ui-state}\\ \hline
\textbf{状态} & \textbf{展示要求} & \textbf{交互要求} \\ \hline
加载中 & 显示骨架屏或明确加载指示 & 防止重复提交产生副作用的操作 \\ \hline
空数据 & 说明当前无数据的原因 & 在合适场景提供创建、搜索或返回入口 \\ \hline
表单错误 & 定位到具体字段并显示可理解文字 & 保留其他合法输入，不清空整张表单 \\ \hline
权限不足 & 显示无权访问，不展示敏感数据 & 引导返回当前角色首页 \\ \hline
状态冲突 & 显示数据已变更或操作已被他人完成 & 刷新最新状态，不覆盖已成功结果 \\ \hline
服务失败 & 显示简洁错误与重试入口 & 写操作仅在存在幂等保护时允许自动重试 \\ \hline
\end{longtable}

\subsection{REST 接口约束}

\begin{longtable}{|L{0.22\linewidth}|L{0.68\linewidth}|}
\caption{REST 接口约束}\label{tab:rest-constraints}\\ \hline
\textbf{项目} & \textbf{要求} \\ \hline
路径 & 新接口使用 \texttt{/api/v1} 版本前缀与资源名词；旧 \texttt{/api} 接口在迁移期保持兼容，不要求一次性改完全部路径。 \\ \hline
HTTP 方法 & GET 查询，POST 创建，PATCH 局部修改，DELETE 删除/逻辑删除；支付、取消、接单等动作可使用动作子资源。 \\ \hline
状态码 & 200/201 表示成功，400 表示参数错误，401 表示未认证，403 表示无权，404 表示资源不存在，409 表示状态/并发冲突，500 表示未处理服务端错误。 \\ \hline
统一响应 & 包含业务码、消息、数据和请求标识；不向客户端暴露堆栈和数据库错误。 \\ \hline
分页 & 页码从 1 开始，默认每页 20 条，单页最多 100 条；响应包含总数、页码、页大小和列表。 \\ \hline
认证与归属 & 需要登录的接口从认证上下文获取用户身份，不信任客户端传入的用户 ID；角色校验后继续校验资源归属。 \\ \hline
\end{longtable}

\subsection{外部服务与降级}

\srsfigure{figures/srs/17-external-services.pdf}{外部服务依赖与降级图}{fig:external-services}

AI、图片识别、语音识别、地图和对象存储必须由后端适配层封装。自动化测试使用可控的 Mock/Stub；任一外部服务失败时，不得导致登录、普通搜索、购物车、订单或文本地址功能不可用。

\section{非功能性需求}

\subsection{安全性与隐私}

\begin{enumerate}
  \item 密码使用 BCrypt 或同等级不可逆安全哈希保存，日志和响应不记录明文密码、完整令牌和外部服务密钥。
  \item 后端对顾客、商家、骑手和管理员执行角色权限及资源归属双重校验；前端隐藏按钮不作为安全保障。
  \item 骑手只在执行活动任务期间获取必要的取餐与收货信息，任务结束、取消或改派后收回访问权。
  \item 上传文件校验媒体类型、扩展名和大小，不使用用户原始文件名作为服务端存储路径。
  \item 对外部 AI 发送数据时实施最小化，不发送密码、完整令牌和与当前请求无关的个人信息。
\end{enumerate}

\subsection{性能与容量}

\begin{longtable}{|L{0.27\linewidth}|L{0.33\linewidth}|L{0.30\linewidth}|}
\caption{性能验收基线}\label{tab:performance}\\ \hline
\textbf{对象} & \textbf{指标} & \textbf{测量条件} \\ \hline
普通查询/写入接口 & 95\% 请求在 2 秒内返回 & 本地验收环境、50 个并发用户、不包含外部 AI/地图调用 \\ \hline
列表接口 & 默认 20 条/页，最多 100 条/页 & 使用初始化演示数据和至少 1000 条订单的测试数据 \\ \hline
AI 客服 & 10 秒超时并进入降级 & 使用可控延迟的测试适配器 \\ \hline
并发骑手接单 & 同一任务恰好一个请求成功 & 至少 20 个并发接单请求 \\ \hline
\end{longtable}

\subsection{数据一致性与可靠性}

\begin{enumerate}
  \item 订单创建、库存预占、优惠券锁定和资产处理使用事务或可恢复的一致性方案，任一关键步骤失败不产生半完成订单。
  \item 支付、取消、退回资产、积分发放、确认收货和骑手接单使用业务唯一键、唯一约束、条件更新或版本号防止重复副作用。
  \item 订单、配送和资产关键变更保存原状态、新状态、执行者、时间和业务原因，不以直接覆盖替代流水。
\end{enumerate}

\subsection{可用性与降级}

外部 AI、图片识别、语音识别、地图或对象存储中任一服务超时/失败时，系统应在等待上限内结束请求并显示降级入口。普通登录、搜索、购物车、订单、文本地址和状态管理不得因此不可用。

\subsection{易用性与可访问性}

\begin{enumerate}
  \item 关键按钮具有默认、禁用、处理中、成功和失败状态，支付、取消和接单等操作防止连续点击。
  \item 店铺休息、商品售罄、订单取消、配送异常等状态同时使用文字和视觉标记，不只依赖颜色区分。
  \item 顾客和骑手核心流程在 375 px 及以上视口宽度下可完成，按钮和输入框不相互遮挡。
\end{enumerate}

\subsection{兼容性}

系统支持验收日期前已发布的 Chrome 和 Edge 最新两个主版本。桌面端在 1366×768 及以上分辨率完成四端核心操作，顾客/骑手端同时满足 375 px 移动视口。

\subsection{可测试性与可维护性}

\begin{enumerate}
  \item P0/P1 需求在实现前明确至少一个正常、一个异常和一个边界测试；测试包含对返回值和数据副作用的明确断言。
  \item 核心 Service 层自动化测试行覆盖率的项目组质量目标为不低于 80\%；覆盖率不替代对业务规则、权限、并发和回退结果的断言。
  \item 外部服务依赖通过适配接口隔离，自动化测试不强制调用真实付费服务。
  \item 需求变更时先修改需求与验收标准，增加会失败的测试，再修改实现并运行受影响的回归测试。
\end{enumerate}

\subsection{AI 安全与可解释性}

AI 生成的商品候选必须同时返回来自业务数据的商品标识、当前价格和推荐理由。工具调用在后端重新执行认证、授权与资源归属校验，不信任 AI 输出的用户 ID 或订单 ID。对话日志不保存密码、令牌和完整收货信息，用户可删除本人历史对话。

\section{验收与需求追踪}

\subsection{核心端到端验收场景}

\begin{longtable}{|L{0.13\linewidth}|L{0.25\linewidth}|L{0.39\linewidth}|L{0.13\linewidth}|}
\caption{核心验收场景}\label{tab:e2e-acceptance}\\ \hline
\textbf{场景} & \textbf{起点} & \textbf{必须完成的结果} & \textbf{级别} \\ \hline
E2E-01 & 顾客浏览店铺 & 加购、试算、创建、支付、商家履约、骑手配送、收货和评价完整闭环 & P0/P1 \\ \hline
E2E-02 & 待支付订单 & 超时或主动取消，库存与优惠券释放，不发生扣款 & P0 \\ \hline
E2E-03 & 已支付待接单订单 & 商家拒单或顾客取消，库存恢复且资产原路返回 & P0 \\ \hline
E2E-04 & 两名上线骑手查看同一任务 & 并发接单恰好一人成功，另一人获得冲突结果 & P0 \\ \hline
E2E-05 & 骑手配送中上报地址异常 & 管理员处理并留存全部轨迹，未授权角色无法改状态 & P0 \\ \hline
E2E-06 & 顾客向 AI 说明预算和忌口 & 返回真实可售商品卡片，经顾客确认后加入普通购物车 & P1 \\ \hline
E2E-07 & AI 或地图服务超时 & 显示降级提示，普通搜索、订单和文本地址仍可使用 & P0 \\ \hline
E2E-08 & 顾客 A/商家 A/骑手 A 猜测他人资源 ID & 读取或写入请求被拒绝且不暴露私有数据 & P0 \\ \hline
\end{longtable}

\subsection{需求追踪矩阵}

\begin{longtable}{|L{0.17\linewidth}|L{0.20\linewidth}|L{0.18\linewidth}|L{0.20\linewidth}|L{0.15\linewidth}|}
\caption{需求追踪矩阵}\label{tab:traceability}\\ \hline
\textbf{需求组} & \textbf{主要页面} & \textbf{主要接口资源} & \textbf{主要数据对象} & \textbf{验收场景} \\ \hline
UC-COM & 登录、个人中心、通知 & users, auth, addresses, notifications & User, Role, Address, Notification & E2E-08 \\ \hline
UC-CUS & 店铺、购物车、订单、评价 & stores, products, carts, orders, reviews & Store, Product, Cart, Order, Review & E2E-01--03 \\ \hline
UC-MER & 店铺、商品、订单工作台 & merchant-applications, stores, products, orders & Merchant\-Profile, Store, Product, Order & E2E-01--03 \\ \hline
UC-RID & 可接任务、当前任务、历史 & rider-applications, delivery-tasks & Rider\-Profile, Delivery\-Task, Delivery\-Exception & E2E-01, 04--05 \\ \hline
UC-ADM & 审核、治理、看板 & audits, users, reviews, delivery-exceptions, statistics & AuditLog, Review, DeliveryException & E2E-05, 08 \\ \hline
UC-AI & AI 客服与智能点餐 & ai-sessions, ai-messages, menu-search, order-query & AiSession, AiMessage, ToolCallLog & E2E-06--07 \\ \hline
UC-MKT/PER & 资产中心、会员、偏好 & wallets, points, coupons, memberships, preferences & Account, Transaction, Coupon, Membership, Preference & E2E-01--03 \\ \hline
\end{longtable}

具体接口路径、数据表名、测试类/方法和最终结果在接口文档、数据库文档和测试报告中按本表需求编号补全，不在 SRS 中绑定具体 Java 类名。

\subsection{单项需求完成标准}

一项 P0/P1 需求仅在同时满足以下条件时标记为完成：

\begin{enumerate}
  \item SRS 与业务规则已确认，需求编号稳定。
  \item 接口行为、权限、输入、响应和错误场景已同步。
  \item 至少有正常、异常和边界自动化测试，关键副作用存在断言。
  \item 前后端联调通过，加载、空、失败和无权状态可正确展示。
  \item 受影响的旧功能回归测试通过，代码、测试和文档已持续提交。
  \item 主责成员能够说明需求、核心规则、实现方式和测试结果。
\end{enumerate}

\section{附录}

\subsection{基线业务配置}

\begin{longtable}{|L{0.28\linewidth}|L{0.20\linewidth}|L{0.42\linewidth}|}
\caption{基线业务配置}\label{tab:baseline-config}\\ \hline
\textbf{配置项} & \textbf{基线值} & \textbf{说明} \\ \hline
手机号 & 11 位数字 & 作为登录账号时全局唯一 \\ \hline
密码长度 & 8--32 个字符 & 至少包含字母和数字，只保存安全哈希 \\ \hline
待支付超时 & 15 分钟 & 超时后取消并释放预占库存/优惠券 \\ \hline
送达后自动完成 & 24 小时 & 顾客未主动确认时由系统完成 \\ \hline
顾客自助取消 & 待支付、待商家接单 & 商家接单进入制作中后不提供自助取消 \\ \hline
评价评分/文字/图片 & 1--5 分；0--500 字；0--3 张 & 图片仅支持 JPG/PNG/WebP，单张最大 5 MB \\ \hline
会员 & 30 天；95\% 折扣 & 只对商品小计打折，不对配送费打折 \\ \hline
优惠券 & 每单最多 1 张 & 必须满足有效期、店铺/平台范围和最低金额 \\ \hline
积分抵扣 & 100 积分 = 1 元；最多 20\% & 上限基于会员和优惠券计算后金额 \\ \hline
积分发放 & 实付每满 1 元发 1 分 & 仅已完成订单发放，重复确认不重复发放 \\ \hline
分页 & 默认 20，最大 100 & 页码从 1 开始 \\ \hline
AI 超时 & 10 秒 & 超时返回降级结果，不自动修改购物车 \\ \hline
\end{longtable}

上述值可通过统一后端配置调整，但每次基线验收必须在测试报告中记录实际使用值，不得由前端页面各自写死不同规则。

\subsection{业务错误类型}

\begin{longtable}{|L{0.21\linewidth}|L{0.22\linewidth}|L{0.47\linewidth}|}
\caption{业务错误类型}\label{tab:error-catalog}\\ \hline
\textbf{类型} & \textbf{HTTP 状态} & \textbf{示例} \\ \hline
参数错误 & 400 & 非法手机号、价格小于等于零、评分越界、日期范围错误 \\ \hline
未认证 & 401 & 未提供凭证、凭证过期或签名无效 \\ \hline
无权限 & 403 & 非管理员审核、非本店商家操作、非本任务骑手更新状态 \\ \hline
资源不存在 & 404 & 订单、商品、地址或配送任务不存在 \\ \hline
业务/并发冲突 & 409 & 重复注册、重复支付、非法状态转换、配送任务已被他人接取 \\ \hline
外部服务不可用 & 503 & AI、地图或对象存储不可用；响应同时提供降级建议 \\ \hline
\end{longtable}

\subsection{配套交付物}

SRS 后续应与数据库设计、REST 接口文档、自动化测试与覆盖率报告、需求变更记录、AI 辅助使用记录、部署说明和交叉验收记录保持需求编号一致。

\section{版本修订记录}

\begin{longtable}{|L{0.16\linewidth}|L{0.14\linewidth}|L{0.46\linewidth}|L{0.14\linewidth}|}
\caption{版本修订记录}\label{tab:revision}\\ \hline
\textbf{日期} & \textbf{作者} & \textbf{修订内容} & \textbf{版本} \\ \hline
2026-09-02 & 项目组 & 形成覆盖四类角色、智能服务和质量要求的第一版文字稿。 & 1.0 \\ \hline
2026-09-02 & 项目组 & 重构协作目录，收敛 P0/P1/P2 范围，明确订单、配送、资产和 AI 安全验收规则。 & 1.1 \\ \hline
\end{longtable}



	% \clearpage

\end{document}                                 % 结束全文
